% Transcribed from CMM127 Midterm.docx.
\IfFileExists{Book\space NL/template/kaobook.cls}{
  \IfFileExists{main.tex}{\documentclass[main.tex]{subfiles}}{
    \IfFileExists{../../main.tex}{\documentclass[../../main.tex]{subfiles}}{
      \IfFileExists{../main.tex}{\documentclass[../main.tex]{subfiles}}{
        \documentclass[Book\space NL/main.tex]{subfiles}
      }
    }
  }
}{
  \IfFileExists{../../template/kaobook.cls}{\documentclass[../../main.tex]{subfiles}}{\documentclass{book}}
}
\IfFileExists{Book\space NL/template/kaobook.cls}{}{\IfFileExists{../../template/kaobook.cls}{}{\InputIfFileExists{review-notes-standalone.tex}{}{\InputIfFileExists{Book\space NL/appendices/appendix 4-Review Notes/review-notes-standalone.tex}{}{}}}}
\begin{document}
\ifSubfilesClassLoaded{
  \ifcsname appendixchapterparent\endcsname
  \else
    \standaloneappendixchaptersetup{D}{2}
  \fi
}{}
\IfFileExists{\booknlroot appendices/appendix 4-Review Notes/images/d2-figure-research-study-design-framework.png}{
  \edef\reviewnotesimagepath{\booknlroot appendices/appendix 4-Review Notes/images/}
}{
  \edef\reviewnotesimagepath{\booknlroot appendices/appendix 4-Review Notes/images/}
}
\chapter{Research Methods and Critical Appraisal}
\labch{appendix-research-methods-critical-appraisal}

\margintoc
\noindent\textit{Research Methods Summary}\hfill\textit{Tai’s Edition}\par\medskip
\section{Introduction to Epidemiology and Research}
\subsection{Research}
\begin{itemize}
\item There are two main types of questions in research: foreground and background questions.
\item Background questions are about aspects of healthcare. For example, questions regarding diagnosis or treatment of a disorder. They contain a question root and an aspect of the condition in question.
\item Foreground questions ask about specific knowledge used in clinical decisions and treatment. This includes the PICO framework: patient, intervention, comparison and outcome. They often lead to research questions.
\item Clinical questions often inform research questions.
\item Research questions consist of a determinant (the thing to be tested), outcome (the thing to be measured) and domain (who the question applies to).
\item The FINER research criteria sets out the aspects of good research questions. Good research questions should be feasible, interesting, novel, ethical and relevant.
\begin{itemize}
\item Feasible: adequate number of subjects, adequate technical expertise, affordable in time and money, manageable in scope and fundable.
\item Interesting: getting the answer intrigues the investigator and her colleagues.
\item Novel: provide new findings, confirm, refute, or extend previous findings, may lead to innovations in concepts of health and disease, medical practice, or methodologies for research.
\item Ethical: a study that the institutional review board will approve.
\item Relevant: likely to have significant impacts on scientific knowledge, clinical practice, or health policy, may influence directions of future research.
\end{itemize}
\end{itemize}
\subsection{Clinical Epidemiology}
\begin{itemize}
\item Epidemiology studies the occurrence of illness or disease, and its relationship to various determinants.
\item Clinical epidemiology is a focused area of epidemiology relevant to clinical practice and patient care. These include diagnosis, etiology, prognosis and therapy.
\item There are two main types of studies: observational and interventional.
\item There are four main types of knowledge:
\begin{itemize}
\item Diagnostic - identify illness based on the clinical profile.
\item Etiologic - explanation of the factors of disease.
\item Prognostic - identify the course and outcomes of a disease.
\item Therapeutic - identify the course and outcomes of a disease with treatment.
\end{itemize}
\end{itemize}
\subsection{Aspects of Research}
\subsubsection{Theoretical Design}
\begin{center}\includegraphics[width=0.82\textwidth,height=0.68\textheight,keepaspectratio]{\reviewnotesimagepath d2-figure-research-study-design-framework.png}\par\captionof{figure}{The design pathway translates a research question into a target population, intended sample, study plan, and measured findings.}\labfig{appendix-figure-d2-research-study-design-framework}\end{center}
\begin{itemize}
\item The theoretical design aspect of research identifies the relations between determinants (things which cause outcomes) that are being investigated in the study. Outcome (y) = f (determinant; x) (see Figure~\ref{fig:appendix-figure-d2-research-study-design-framework}).
\begin{itemize}
\item Causal research identifies the potential determinants which lead to an outcome (cause and effect). In this type of research, confounders (factors that influence results) are an issue. Confounders are also known as extraneous determinants (ED). Causal research is often used to treat.
\item Descriptive research studies the factors associated with a disease which may or may not be directly causal. This type of research is used in the diagnosis and prognosis of disease (see Figure~\ref{fig:appendix-figure-d2-research-causal-and-descriptive-designs}).
\end{itemize}
\end{itemize}
\begin{center}\includegraphics[width=0.82\textwidth,height=0.68\textheight,keepaspectratio]{\reviewnotesimagepath d2-figure-research-causal-descriptive-designs.png}\par\captionof{figure}{Causal designs isolate an exposure–outcome relationship, whereas descriptive designs consider several determinants associated with diagnosis or prognosis.}\labfig{appendix-figure-d2-research-causal-and-descriptive-designs}\end{center}
\subsubsection{Domain and Sample}
\begin{itemize}
\item The domain of a study is all the patients to which a study is applicable.
\item The inclusion criteria specify the population relevant to the research question and eligible for the study. Characteristics include:
\begin{itemize}
\item Demographic characteristics
\item Clinical characteristics
\item Geographic characteristics
\item Temporal characteristics
\end{itemize}
\item The exclusion criteria specify which subset of the inclusion criteria will not be included in the study. This may be due to:
\begin{itemize}
\item A high likelihood of being lost to follow-up
\item An inability to provide good data
\item  Being at high risk of possible adverse effects
\end{itemize}
\end{itemize}
\subsubsection{Study Outcomes}
\begin{itemize}
\item Outcomes of a study (endpoints/events) are the variables which are monitored to document the impacts of a treatment on a population. Typical examples of outcomes are cure, clinical worsening (deterioration), and mortality.
\item There are four main types of outcomes:
\begin{itemize}
\item Primary outcomes are outcomes directly related to the research question. Generally, primary outcomes should be patient centered.
\item Secondary outcomes are outcomes which are monitored to assist in interpretation of primary outcomes.
\item Surrogate outcomes are biomarkers intended to substitute for a clinical outcome.
\item Composite outcomes are made up of multiple variables.
\end{itemize}
\item Outcomes should be chosen before looking at the results (a priori) to avoid drawing false conclusions through the Texas sharpshooter fallacy.
\item The sample size calculation should be carried out to detect a clinically relevant effect of the intervention on the primary outcome.
\item Calculations can also be made for secondary outcome variables, which may increase the sample size but also increase trial validity.
\item The outcome chosen should be based on the research question and hypothesis. If the outcome is clinical prognosis, do not choose teeth colour.
\end{itemize}
\subsubsection{Data Collection Design}
\begin{itemize}
\item The data collection design for research depends on a number of factors, including domain, time constraints and budget.
\item Data collection methods are classified by \textbf{time}, \textbf{study population} (sampling/census) and \textbf{type} (observational/experimental).
\item Cross-sectional studies investigate subjects and outcomes simultaneously at one point in time. They associate one variable with another.
\end{itemize}
\begin{center}\small
\renewcommand{\arraystretch}{1.15}
\captionof{table}{Strengths and limitations of a one-time-point population assessment.}\labtab{appendix-table-d2-cross-sectional-study-strengths-limitations}
\begin{tabular}{p{0.45\linewidth}p{0.45\linewidth}}\toprule
\textbf{Advantages of cross-sectional studies} & \textbf{Disadvantages of cross-sectional studies} \\\midrule
\textbullet\ Relatively inexpensive and takes little time to conduct.\par
\textbullet\ Can estimate prevalence of the outcome because the sample is usually taken from the whole population.\par
\textbullet\ Many outcomes and risk factors can be assessed.\par
\textbullet\ Useful for public-health planning, understanding disease aetiology, and generating hypotheses.\par
\textbullet\ There is no loss to follow-up. & \textbullet\ Difficult to make causal inference.\par
\textbullet\ Only a snapshot; results may differ if another time frame had been chosen.\par
\textbullet\ Prevalence--incidence (Neyman) bias may under-represent people with longer-lasting disease when a risk factor also causes death. \\\bottomrule
\end{tabular}
\end{center}
\begin{itemize}
\item If a study is not cross-sectional, it is either prospective or retrospective.
\item Prospective studies investigate data which is collected after formation of the study.
\item Retrospective studies use data already collected before the study formation.
\item Sampling can either be done by census (all members of a population) or sampling (some members of a population).
\end{itemize}
\begin{center}\includegraphics[width=0.82\textwidth,height=0.68\textheight,keepaspectratio]{\reviewnotesimagepath d2-figure-study-design-retrospective-prospective-case.png}\par\captionof{figure}{The timeline distinguishes exposure-first cohort follow-up from case--control sampling, which begins with outcome status and looks back for exposure.}\labfig{appendix-figure-d2-study-design-retrospective-prospective-case-control-timelines}\end{center}
\begin{itemize}
\item The study types outlined above are observational. However, there are also studies where patients are given an intervention. These studies are called experimental studies (see Figure~\ref{fig:appendix-figure-d2-random-allocation-study-groups-outcomes}).
\item Experimental studies assign an intervention; randomized controlled trials use random allocation, while quasi-experimental designs do not.
\end{itemize}
\begin{center}\includegraphics[width=0.82\textwidth,height=0.68\textheight,keepaspectratio]{\reviewnotesimagepath d2-figure-research-study-validity-inference-pathway.png}\par\captionof{figure}{Two inferential steps separate a research question from observed findings: designing the study and implementing that plan.}\labfig{appendix-figure-d2-research-study-validity-inference-pathway}\end{center}
\begin{itemize}
\item Internal validity concerns whether a study's results are valid for its participants, with minimal bias from design, conduct, measurement, or analysis (see Figure~\ref{fig:appendix-figure-d2-research-study-validity-inference-pathway}).
\item External validity refers to the study design’s ability to reflect the real world (see Figure~\ref{fig:appendix-figure-d2-research-study-validity-inference-pathway}).
\item There are two types of error: random error (which reduces precision) and systematic error (which reduces accuracy).
\item Systematic error may either be fixed amounts or proportions.
\item Random error arises by chance and may distort results in either direction. Larger samples can improve precision and statistical power, but do not eliminate random error.
\item Systematic errors are errors in the study design which distort study findings in one direction. Systematic errors include biases such as selection bias, information bias and confounder bias. To eliminate systematic error, study accuracy should be increased by redesign to increase conclusion validity.
\item Precision refers to the repeatability of a value, while accuracy refers to distance from the true value.
\item Precision is evaluated by comparison among repetitions. Accuracy is determined by comparison with a “gold standard.”
\item Both types of error may occur from the observer, subject and instrument.
\end{itemize}
\begin{itemize}
\item The generalizability of a study refers to applicability in the real world, relevance.
\item Generalizability is a product of the study design and research question, along with accuracy and precision.
\item Research questions should be clinically relevant and directly lead to the study design.
\item Study designs should have relevant patient domains to the target population and good monitored outcomes and Relevant predictors/determinants/interventions.
\item To be generalizable, a study needs clinical relevance, validity and precision.
\end{itemize}
\begin{itemize}
\item The DEPTh (Diagnostic, Etiologic, Prognostic, Therapeutic) model outlines the types of studies which are typically done in clinical research.
\end{itemize}
\section{Critical Appraisal}
\begin{itemize}
\item Evidence based medicine refers to the application of the best available research to clinical care. It integrates evidence with clinical expertise and patient values.
\end{itemize}
\begin{center}\includegraphics[width=0.82\textwidth,height=0.68\textheight,keepaspectratio]{\reviewnotesimagepath d2-figure-evidence-based-medicine-clinical-expertise.png}\par\captionof{figure}{Evidence-based decisions balance the best available research with practitioner expertise and the patient's circumstances and preferences.}\labfig{appendix-figure-d2-evidence-based-medicine-clinical-expertise-evidence-values}\end{center}
\begin{itemize}
\item The application of evidence based medicine is based on the five As:
\begin{itemize}
\item Ask: Converting the need for information into an answerable question.
\item Acquire: Tracking down the best evidence with which to answer that question.
\item Appraise: Critically appraising that evidence for its validity (closeness to the truth), impact (size of the effect), and applicability (usefulness in our clinical practice).
\item Apply: Integrating the critical appraisal with our clinical expertise and with our patient's unique biology, values, and circumstances.
\item Assess: Evaluating our effectiveness and efficiency in executing steps 1 to 4 and seeking ways to improve them both for next time.
\end{itemize}
\item \textbf{Ask} (in the 5As): see background, foreground questions and PICO. Also note that background questions are mostly asked in learning about a disease, and foreground questions used to treat.
\item \textbf{Acquire} the best evidence using the hierarchy of evidence (see Figures~\ref{fig:appendix-figure-d2-evidence-hierarchy-pyramid} and~\ref{fig:appendix-figure-d2-evidence-hierarchy-search-and-processing}).
\end{itemize}
\begin{center}\includegraphics[width=0.82\textwidth,height=0.68\textheight,keepaspectratio]{\reviewnotesimagepath d2-figure-evidence-hierarchy-pyramid.png}\par\captionof{figure}{The pyramid summarizes a commonly taught ordering of evidence sources; the most appropriate design depends on the clinical question.}\labfig{appendix-figure-d2-evidence-hierarchy-pyramid}\end{center}
\begin{center}\includegraphics[width=0.82\textwidth,height=0.68\textheight,keepaspectratio]{\reviewnotesimagepath d2-figure-evidence-hierarchy-search-processing.png}\par\captionof{figure}{The diagram links question-specific evidence hierarchies to primary studies, evidence synthesis, and progressively processed research resources.}\labfig{appendix-figure-d2-evidence-hierarchy-search-and-processing}\end{center}
\noindent\textit{Source attribution reproduced from the original diagram: User's Guide to the Medical Literature, 3rd edition.}
\begin{itemize}
\item Non-appraised (primary) research is research which is unfiltered and is often a source of data for secondary (pre-appraised) research.
\item Secondary research is often more information dense and verified, but at the risk of bias. Examples are seen above.
\item Searching for a question in literature involves extracting keywords, creating syntax (with appropriate boolean operators: AND, OR, NOT) and searching.
\item After finding results, the quality of literature must be ascertained.
\item Factors Determining Journal Quality
\begin{itemize}
\item Peer Review Process – Ensures the credibility and accuracy of published research.
\item Journal Databases – Inclusion in recognized databases indicates reliability and impact:
\item[]
\begin{itemize}
\item ISI (Institute for Scientific Information®)
\item Scopus – The largest abstract and citation database of peer-reviewed literature.
\item TCI (Thai Journal Citation Index Centre)
\end{itemize}
\item Journal Impact Factor (JIF) – A journal-level citation metric based on the frequency with which its articles are cited during a defined period; it is not a direct measure of the quality of an individual article.
\item H-Index – A metric that reflects both the productivity and citation impact of a researcher's publications. A scientist has an index h if h of their N papers have at least h citations each, while the remaining (N-h) papers have no more than h citations each.
\item Quartile Ranking (Q1–Q4) – Journals may be grouped into quartiles based on citation metrics within a category. Peer review, indexing, and journal-level metrics can inform an appraisal but do not establish that an individual article is reliable or methodologically sound.
\end{itemize}
\item Appraisal of primary research aims to assess validity (distance from the truth), impact (size of effect), and applicability.
\item To effectively appraise, ask yourself “Why am I reading?” Stay focused and make sure that you get the information that you are looking for. You should also be selective of the paper chosen.
\item A research article is composed of:
\begin{itemize}
\item A title.
\item An abstract.
\item Background / introduction.
\item Materials and methods.
\item Results.
\item Discussion.
\item Conclusion.
\end{itemize}
\end{itemize}
\subsection{The literature review process}
\begin{itemize}
\item[]
\begin{itemize}
\item Be selective of papers chosen to read.
\item Primary literature consists of full papers such as original or research papers. They are most up-to-date but contain only one direction of thought. Many should be read to get the whole picture.
\item Secondary literature is a summary of primary literature such as literature reviews and textbooks. They may contain bias, but are often a good place to start (not used for further research though).
\end{itemize}
\item Critical reading is a technique used to identify ideas in a text, while critical thinking evaluates the merit of information and ideas.
\item In critical appraisal, a fact level of understanding and an interpretation level of understanding can be achieved. It is the difference between remembering the text and understanding it and discussing aspects of the paper, along with your interpretation.
\item The goal of critical reading is recognizing purpose, tone, persuasive elements and bias in a text.
\item Checklist for reading critically:
\begin{itemize}
\item A specific topic must be addressed.
\item Terms must be clearly defined.
\item Evidence must be presented.
\item Common knowledge must be accounted for.
\item Exceptions must be explained.
\item Causes must be shown to precede effects and to be capable of the effect.
\item Conclusions must be shown to follow logically from earlier arguments and evidence.
\end{itemize}
\end{itemize}
\subsection{Appraising Aspects of a Research Paper}
\subsubsection{Abstract}
\begin{itemize}
\item The abstract should be well written, and contain the aim of the study.
\item The abstract should contain a summarized version of all the information, Including the material, methods, results and conclusion.
\item Study design, grouping, randomization, procedures and measuring tools should all be included, along with measured variables and their significance.
\item The conclusion should answer the question of interest. 
\item There are two main types of abstracts, which are structured and unstructured.
\end{itemize}
\begin{center}\small
\renewcommand{\arraystretch}{1.15}
\captionof{table}{Questions for orienting a first reading of a research paper.}\labtab{appendix-table-d2-research-paper-title-author-source-purpose-abstract}
\begin{tabular}{p{0.18\linewidth}p{0.70\linewidth}}\toprule
\textbf{Title} & What information does the title give about the article's focus, methodology, and subject? What does it lead you to expect? \\\midrule
\textbf{Author(s)} & Who wrote the text? What are the author's discipline, expertise, and background? Names may become familiar as knowledge of a field expands. \\
\textbf{Source} & Where was the work published, and what is its original source? Is it from an academic journal, textbook, or popular magazine? The publication may lend credibility to or detract from the material. \\
\textbf{Purpose} & Why did the author write the paper or conduct the study? Does it refute another interpretation or offer new information? The abstract usually provides clues. \\
\textbf{Abstract} & Read it carefully for an uncritical summary of the paper's subject and content, including the study rationale, main results, how they were obtained, and their interpretation. \\\bottomrule
\end{tabular}
\end{center}
\subsection{Background/Introduction}
\begin{itemize}
\item Discusses earlier works in the area under study, while also identifying gaps in current knowledge.
\item The introduction expresses the rationale and purpose of the study.
\item Indicates whether data was searched for a result, or searched with a prior purpose.
\item Should outline potential benefits and drawbacks of the method or observations.
\end{itemize}
\subsubsection{Materials and Methods}
\begin{itemize}
\item Identifies how the study was conducted in detail.
\item Information on study design, population, sample size and interventions presented.
\item All measurement approaches should be clear.
\end{itemize}
\begin{center}\small
\renewcommand{\arraystretch}{1.15}
\captionof{table}{Methodological questions for assessing whether a study design addresses its objective.}\labtab{appendix-table-d2-critical-appraisal-methods-prompts}
\begin{tabular}{p{0.20\linewidth}p{0.68\linewidth}}\toprule
\textbf{Central objective} & What is the central objective of the paper, and what is the major question being addressed? \\\midrule
\textbf{Hypothesis and follow-through} & Does the hypothesis address the main study question? Are there underlying assumptions, and has it been tested properly? \\
\textbf{Use of controls} & What controls were used in the study? Are they adequate? Why or why not? \\
\textbf{Procedures} & What procedures and methodologies were used to address the central question? What experimental procedures were conducted, and what is the rationale behind each? Has each procedure been explained properly and does it test the central hypothesis? What limitations do the procedures or methodologies have? \\\bottomrule
\end{tabular}
\end{center}
\subsubsection{Results}
\begin{itemize}
\item Indicates the endpoints of the subjects.
\item May contain raw data and explanations of statistical analysis.
\item Results are often shown in relevant tables, diagrams and graphs.
\end{itemize}
\begin{center}\small
\renewcommand{\arraystretch}{1.15}
\captionof{table}{Prompts for checking the clarity, consistency, and evidential value of reported results.}\labtab{appendix-table-d2-critical-appraisal-results-and-display-prompts}
\begin{tabular}{p{0.20\linewidth}p{0.68\linewidth}}\toprule
\textbf{Results} & Do the results adequately and accurately describe the data? Are they consistent and do they fit the context of the paper? Are they sufficient to test the central hypothesis and allow the researcher or reader to accept or reject it? Do they substantiate the author's claims? Why or why not? \\\midrule
\textbf{Figures and tables} & Are the figures and tables helpful to understanding and ease of reading, or are they redundant? Are they organized effectively? \\\bottomrule
\end{tabular}
\end{center}
\subsubsection{Discussion}
\begin{itemize}
\item Includes comparison between established data and new discoveries. Also discusses clinical relevance.
\item Discusses potential limitations and potential or necessitation for further studies.
\end{itemize}
\begin{center}\small
\renewcommand{\arraystretch}{1.15}
\captionof{table}{Questions for judging interpretation and the overall usefulness of an article.}\labtab{appendix-table-d2-critical-appraisal-discussion-summary-prompts}
\begin{tabular}{p{0.20\linewidth}p{0.68\linewidth}}\toprule
\textbf{Discussion} & Are the author's conclusions logical based on the results, or do they seem too optimistic, large, or unjustified? Are facts and opinions clearly separated, or are they difficult to distinguish? \\\midrule
\textbf{Summary and conclusion} & Is the article useful for your purposes? Does it fit your knowledge or current understanding of the topic, or is it something new? If it is new, does it help you understand the topic, or do you need more information? \\\bottomrule
\end{tabular}
\end{center}
\subsubsection{Evaluation}
\begin{itemize}
\item Analyze the research question and specificity of the effect studied. Could it be caused by something else?
\item Look at and determine if the methods used in the study were appropriate and the best approach considering limited resources.
\item Look at the objectiveness of the data interpretation, and whether conclusions were justified.
\item Were the original questions answered?
\item Implications of the study on practice in the field.
\item Potential suggestions for further research as indicated by the paper.
\end{itemize}
\subsection{Research misconduct}
\begin{itemize}
\item Defined as actions and behaviors by researchers that fail to honor the integrity of research.
\item \textbf{Fabrication}: Making up data without conducting the required research.
\item \textbf{Falsification}: Manipulating data, materials, processes, or equipment to reach a predefined conclusion.
\item \textbf{Plagiarism}: Using another person’s ideas, content, writing, processes, or results without proper credit (includes self-plagiarism).
\item \textbf{Authorship}:
\begin{itemize}
\item Assigning false authorships without adequate contribution.
\item Mentioning authors without their consent.
\item Omitting original contributors.
\item Incorrectly ordering or naming authors.
\end{itemize}
\item \textbf{Conflicts of interest}: Financial, personal, or professional conflicts affecting research integrity.
\item \textbf{Approvals}: Non-compliance with ethical mandates related to human or animal research.
\end{itemize}
\subsection{Steps to Critical Appraisal}
\subsubsection{The Research Question}
\begin{itemize}
\item See PICO above.
\end{itemize}
\subsubsection{The Study Method}
\begin{center}\includegraphics[width=0.82\textwidth,height=0.68\textheight,keepaspectratio]{\reviewnotesimagepath d2-figure-study-validity-external-internal-inference.png}\par\captionof{figure}{External validity concerns the inference from the target population to the study plan; internal validity concerns the inference from the plan to the observed study findings.}\labfig{appendix-figure-d2-study-validity-external-internal-inference}\end{center}
\begin{center}\small
\renewcommand{\arraystretch}{1.15}
\captionof{table}{Design and analysis strategies for considering random error, bias, and variable selection.}\labtab{appendix-table-d2-research-error-mitigation-by-phase}
\begin{tabular}{p{0.18\linewidth}p{0.31\linewidth}p{0.31\linewidth}}\toprule
\textbf{Potential error} & \textbf{Design phase} & \textbf{Analysis phase} \\\midrule
Chance (random error) & Increase sample size and use other strategies to increase precision. & Calculate p-values and confidence intervals, and interpret them in the context of prior evidence. \\
Bias (systematic error) & Consider the potential consequences of differences between the research question and study plan; alter the plan if necessary. & Check consistency with other studies, especially those using different designs. \\
Assessing possible bias & Collect additional data to assess the extent of possible biases. & Analyze additional data to determine whether potential biases occurred. \\
Variable selection and adjustment & Do not use variables affected by the predictor of interest as inclusion criteria or matching variables. & Do not control for variables affected by the predictor variable. \\\bottomrule
\end{tabular}
\end{center}
\paragraph{Common Biases}
\begin{itemize}
\item Selection bias. In essence, the sample selected for study is not representative. Examples include ascertainment bias, sampling bias, competing risks bias, volunteer bias, nonresponse bias, loss to follow-up bias, attrition bias, prevalence-incidence (Neyman) bias, survivor treatment selection bias and overmatching bias.
\item Information bias refers to inaccuracies which may present in the information collected from a study. Examples include: recall bias, interviewer bias, observer bias, lead-time bias, length-time bias, diagnostic bias, Will Rogers phenomenon, family information bias.
\item Intervention biases are differences due to the way an intervention is carried out. They include: allocation bias, compliance bias, proficiency bias, contamination bias.
\item Confounding occurs when a third variable is associated with both the exposure and outcome, is not an intermediate on the causal pathway, and distorts their observed association if not appropriately addressed.
\item Attrition bias is caused by participant dropouts, which may potentially only show the favourable effects of a treatment.
\item Bias should be prevented at the study design and data handling stages.
\item Study designs are as seen above.
\end{itemize}
\subsubsection{Results of the Study}
\begin{itemize}
\item Statistical significance alone may not indicate clinical significance.
\item Statistical significance describes how incompatible the observed data are with a specified null model; it does not establish causation or indicate the probability that the finding occurred by chance.
\item Clinical significance indicates the effectiveness and utility of findings in real practice.
\item Statistical significance \(\neq\) clinical significance. Small improvements may not be meaningful.
\end{itemize}
\subsubsection{Applicability of Results}
\begin{itemize}
\item Simply asks whether the results of a study are applicable to the patient, condition and other factors at hand.
\end{itemize}
\begin{itemize}
\item JBI's tools for critical appraisal are often used for various study types.
\end{itemize}
\section{Study Designs}
\begin{itemize}
\item There are 4 main study designs:
\begin{itemize}
\item Diagnostic studies describe and predict the probability of a target disease being present based on clinical and non clinical profile.
\item Prognostic research identifies the potential outcomes and predicts the course of a disease from clinical and non clinical profiles.
\item Etiologic research identifies the causes and contributing factors for the occurrence of a disease.
\item Interventional research identifies the effect of a treatment on the course of a disease, while also explaining the cause and effect. It also describes the course of the disease given the various treatment options and clinical or non clinical profile.
\end{itemize}
\end{itemize}
\subsection{Diagnostic Research}
\begin{itemize}
\item Diagnostic research focuses on the diagnosis of disease.
\item The onset of a disease is often asymptomatic; it is known as the preclinical phase.
\item At the clinical (symptomatic) phase, the disease is diagnosed and therapy begins.
\item Screening aims to detect disease or risk markers before symptoms develop; apparent benefit must be distinguished from lead-time and overdiagnosis effects.
\item Screening refers to the searching for an unrecognized (asymptomatic) defect or disease through rapid testing of everyone meeting the preliminary criteria. These tests are not used to treat.
\item An example is mammography for cancer screening.
\item Diagnostic tests are used to diagnose diseases in symptomatic people. They are also used to inform treatments and confirm or refute the diagnosis.
\item The gold standard test is the best test for a diagnosis. However, diagnostic gold standards are often costly and invasive.
\end{itemize}
\begin{center}\small
\renewcommand{\arraystretch}{1.15}
\captionof{table}{Contrasting population screening with tests used to evaluate an individual patient.}\labtab{appendix-table-d2-screening-and-diagnostic-test-comparison}
\begin{tabular}{p{0.45\linewidth}p{0.45\linewidth}}\toprule
\textbf{Screening tests} & \textbf{Diagnostic tests} \\\midrule
Done on apparently healthy people. & Done on those with indications of disease. \\
Applied to groups. & Applied to a single patient. \\
Positive results usually require diagnostic follow-up. & Results inform diagnosis and subsequent clinical decisions. \\
Based on one criterion or cut-off point. & Based on a number of findings. \\
Designed for use in an eligible asymptomatic population. & Selected for a patient's clinical presentation and question. \\
May prompt diagnostic assessment; does not alone establish a diagnosis. & Findings are interpreted with the clinical context to guide diagnosis and care. \\\bottomrule
\end{tabular}
\end{center}
\begin{itemize}
\item The lead time is the time gained by screening when compared to the detection from symptoms (early diagnosis of disease through screening tests).
\item Earlier diagnosis can improve outcomes when effective early treatment changes the disease course; screening benefit must be assessed using meaningful outcomes and accounting for lead-time and overdiagnosis bias.
\item Considerations in the diagnosis of disease include:
\begin{itemize}
\item The importance of disease diagnosis.
\item The natural history of the disease should be well known (the history of disease without treatment).
\item The latent or early symptomatic states of the disease should be detectable.
\item The presence of a diagnostic tool capable of recognizing a disease before onset should be available.
\item Morbidity and mortality should be reduced by early treatment.
\end{itemize}
\item The key considerations for screening tests include:
\begin{itemize}
\item Acceptability for the patient, including low discomfort or embarrassment,
\item Repeatability and reliability. Low inter/intraobserver variation. This means that tests give consistent results when repeated on the same individual or material under the same conditions,
\item Validity and accuracy: High sensitivity and specificity mean that a test measures what it intends to measure and successfully distinguishes between people with and without the disease.
\item Simplicity.
\item Safety.
\item Speed.
\item Pricing.
\end{itemize}
\item A new diagnostic test may be compared against the gold standard by performing both tests on the same patients and comparing outcomes as measured by the test.
\item A standard 2×2 table and sums may be used in the comparison (positive and negative of each test crossed).
\item In diagnostic research, there must be a necessity for the test, and the condition should be treatable. New tests should also have an advantage over the reference standard.
\item The sensitivity of a test may be calculated from the number of measured positives divided by the total number of disease positive patients.
\item The specificity of a test is calculated from the number of measured negative patients over total negative patients.
\item The cutoff point for a test is often described as a compromise between sensitivity and specificity. A more sensitive test is often less specific, while a more specific one is less sensitive.
\item Generally, if the penalty for missing a case is high, a high sensitivity is preferred. Especially to screen the population.
\item If the aim is definite diagnosis, false positives must be minimized. This means high specificity.
\item Validity is the accuracy of a test (how close to the true value it is). Reliability is the repeatability of a test.
\item Positive predictive value (PPV) is the proportion of positive results that are true positives: \(TP/(TP+FP)\).
\item Negative predictive value (NPV) is the proportion of negative results that are true negatives: \(TN/(TN+FN)\).
\item At a fixed threshold, sensitivity and specificity are often less dependent on prevalence than predictive values, but may vary across populations because of differences in disease spectrum and patient characteristics. PPV and NPV vary with prevalence.
\item Changing a test threshold usually trades sensitivity against specificity: a lower threshold tends to increase sensitivity and reduce specificity, while a higher threshold tends to do the opposite.
\end{itemize}
\begin{center}\small
\renewcommand{\arraystretch}{1.15}
\captionof{table}{Diagnostic biases that can distort estimates of test performance.}\labtab{appendix-table-d2-diagnostic-test-biases}
\begin{tabular}{p{0.18\linewidth}p{0.31\linewidth}p{0.31\linewidth}}\toprule
\textbf{Type} & \textbf{Definition} & \textbf{Impact} \\\midrule
Spectrum bias & The study population does not represent the full spectrum of disease severity and normal cases seen in clinical practice. & May overestimate sensitivity and specificity because real-world cases include borderline or atypical presentations. \\
Verification bias (work-up or referral bias) & Only certain patients, often those with positive test results, undergo the gold-standard test. & Sensitivity and specificity may be distorted by selective verification. \\
Observer bias (interpretation bias) & Knowledge of clinical information influences interpretation of a diagnostic test. & Can increase false positives or false negatives; blinding the interpreter to clinical details helps reduce this bias. \\\bottomrule
\end{tabular}
\end{center}
\subsection{Etiologic Research}
\begin{itemize}
\item Etiologic research identifies the potential causes of a disease, or contributing factors.
\item The key goal of etiologic research is to establish causality between two factors. It should indicate that the determinant-outcome relationship is not explained by other factors.
\item Etiologic research may be experimental or observational. However, most are observational due to ethical considerations.
\item A cohort study investigates the end results of two groups; those with exposure and those lacking exposure to a factor. The end incidence for a condition in both groups is then compared.
\begin{itemize}
\item The risk ratio is the incidence risk in the exposed group divided by the incidence risk in the unexposed group. A rate ratio compares incidence rates calculated using person-time in both groups.
\item Cohort studies can be either prospective or retrospective.
\item Retrospective cohort studies require data to have been collected and followed up upon in the past. It is only possible if adequate data has already been collected.
\item Advantages and disadvantages of cohort studies are summarized in Table~\ref{tab:appendix-table-d2-cohort-study-strengths-limitations}.
\begin{center}\small
\renewcommand{\arraystretch}{1.15}
\captionof{table}{Advantages and limitations of following an exposed cohort over time.}\labtab{appendix-table-d2-cohort-study-strengths-limitations}
\begin{tabular}{p{0.45\linewidth}p{0.45\linewidth}}\toprule
\textbf{Advantages of cohort studies} & \textbf{Disadvantages of cohort studies} \\\midrule
\textbullet\ The temporal dimension, whereby exposure is seen to occur before outcome, gives some indication of causality.\par
\textbullet\ Can be used to study more than one outcome.\par
\textbullet\ Good for the study of rare exposures.\par
\textbullet\ Can measure change in exposure and outcome over time.\par
\textbullet\ Incidence of outcome can be measured. & \textbullet\ Costly (less so for retrospective studies) and may take a long time, particularly where outcome onset can occur both early and late in life.\par
\textbullet\ Require accurate records for retrospective studies.\par
\textbullet\ When studying rare outcomes, a very large sample size is required.\par
\textbullet\ Prone to dropout.\par
\textbullet\ Changes in disease aetiology over time may be hard to disentangle from changes observed as age increases.\par
\textbullet\ Selection bias: a difference in outcome incidence between those who participated and those who did not would give biased results. \\\bottomrule
\end{tabular}
\end{center}
\begin{itemize}
\item Selective withdrawal or loss to follow-up may occur.
\item Cohort studies can estimate incidence risks and rates when follow-up and person-time are measured appropriately.
\end{itemize}
\end{itemize}
\item Case--control studies select participants according to outcome status and compare prior or otherwise measured exposures between cases and controls. Exposure data are often collected retrospectively, but this is not required by the design.
\begin{itemize}
\item Two groups are compared; one group with a disease and one without.
\item Case-controlled studies compare the prevalence of risk factors in subjects with and without the disease; case and control.
\item It is often used in the study of rare disease.
\item Because participants are selected by outcome status, a conventional case--control study cannot directly estimate incidence or risk. It commonly uses the odds ratio, comparing exposure odds among cases with exposure odds among controls.
\item Advantages and disadvantages of case--control studies are summarized in Table~\ref{tab:appendix-table-d2-case-control-study-strengths-limitations}.
\begin{center}\small
\renewcommand{\arraystretch}{1.15}
\captionof{table}{Strengths and limitations of selecting participants by outcome status.}\labtab{appendix-table-d2-case-control-study-strengths-limitations}
\begin{tabular}{p{0.45\linewidth}p{0.45\linewidth}}\toprule
\textbf{Advantages of case--control studies} & \textbf{Disadvantages of case--control studies} \\\midrule
Quick and inexpensive; particularly suited to the study of rare diseases because cases are selected at the outset. & \textbullet\ Difficulties in overcoming potential bias and confounding.\par \textbullet\ Selecting representative cases and controls is often difficult.\par \textbullet\ Causality cannot be inferred, and exposure--outcome chronology is unavailable.\par \textbullet\ Inefficient for studying rare risk factors.\par \textbullet\ Often not population-based, so incidence cannot be calculated. \\\bottomrule
\end{tabular}
\end{center}
\end{itemize}
\item Field trials (preventive prophylactic trials) refers to trials in which a healthy population is randomly assigned an intervention (for example vitamins) and studied over time to measure endpoint outcomes.
\end{itemize}
\subsubsection{Biases in Etiologic Research}
\begin{itemize}
\item Selection bias is when errors in the selection procedure for a study affect the representativeness of the results. It also occurs when factors affect study participation.
\item Information bias describes the errors in measurement and collection of information, and subsequent misclassifications regarding the exposure or the disease.
\item Confounders are determinants not in the chain of effect but are capable of influencing the final measured outcome. They must also be associated with the exposure.
\end{itemize}
\subsection{Therapeutic Research}
\begin{itemize}
\item Evaluates the effectiveness of a given treatment. Can alter patient prognosis.
\item Includes an intervention, which may be medical, surgical or behavioral.
\item Helps to indicate the effectiveness and safety of a drug in people.
\item A randomized controlled trial assigns participants to either a treatment or control group randomly. It then observes the difference between the two groups. It is the gold standard of trials.
\item Efficacy refers to how well an intervention works under ideal, controlled conditions, such as in a clinical trial. It measures the maximum potential benefit of treatment when patients strictly follow the protocol, demonstrating a treatment can work.
\item Effectiveness indicates how well a treatment works under real world settings, under influence from patient adherence, comorbidities, and variations in healthcare delivery.
\item Effectiveness indicates if a treatment does work in reality.
\end{itemize}
\begin{itemize}
\item Key elements in a clinical trial are the population, contrast group and measurement of the outcome.
\end{itemize}
\subsubsection{Clinical Trial Phases}
Before clinical trials, preclinical laboratory studies, which may take several years, provide information about dosing and toxicity. The main clinical trial phases are:
\begin{description}
\item[Phase 1.] Primarily evaluates safety, tolerability, and dosing in people; participant numbers and duration vary by intervention and study design.
\item[Phase 2.] Further evaluates safety and dose, and examines preliminary efficacy; participant numbers and duration vary.
\item[Phase 3.] Evaluates efficacy and safety in larger, often multi-site populations; participant numbers and duration vary.
\item[Phase 4 (post-marketing).] Gathers information about safety and effectiveness in real-world populations, including effects across different populations and side effects associated with long-term use.
\end{description}
\begin{center}\includegraphics[width=0.82\textwidth,height=0.68\textheight,keepaspectratio]{\reviewnotesimagepath d2-figure-clinical-trial-development-phases.png}\par\captionof{figure}{Clinical-trial phases progress from initial human safety evaluation to larger efficacy studies and post-marketing surveillance; the sample sizes shown are illustrative.}\labfig{appendix-figure-d2-clinical-trial-development-phases}\end{center}
\subsubsection{Population}
\begin{itemize}
\item The population of a study should be patients currently requiring the treatment, and clear inclusion/exclusion criteria should be set in a manner appropriate to the hypothesis.
\item The sampling strategy should suit the target population and recruitment setting. Consecutive recruitment of eligible patients can reduce selective inclusion within an accessible clinical population, but does not guarantee representativeness of a wider target population. Other strategies include random, stratified, and systematic sampling.
\item Sample size calculation is done by taking the effect size (the difference we expect to detect between groups), variability/standard deviation, significance level/alpha, power/1-beta and study design into account.
\item A Type I error occurs when a true null hypothesis is rejected. The prespecified significance level \(\alpha\) is the long-run probability of this error under the null; the observed p-value is not the probability that a result is a false positive. A threshold such as 0.05 is conventional, not universal.
\item A Type II error occurs when a false null hypothesis is not rejected. Its probability is \(\beta\), and power is \(1-\beta\). A nonsignificant result does not prove that the null hypothesis is true.
\end{itemize}
\subsubsection{Randomization}
\begin{itemize}
\item The generation of a contrast group balances the baseline characteristics of subjects (especially prognostic factors). Minimizes confounders.
\item Randomization acts as the point of departure for the two groups.
\item In an equal-allocation design, each participant has the same probability of assignment to either group. Randomization balances prognostic factors on average across repeated trials, but does not guarantee balance in an individual sample.
\item Randomization is the point of departure between each group.
\item Randomization methods include simple (unrestricted) randomization and restricted randomization.
\item Unrestricted randomization may result in unequal participants in the control and treatment groups. It may also result in different patient populations.
\item Restricted randomization can be achieved through blocked permutation randomization and stratified randomization.
\item Block randomization is where blocks of permutations of a particular size are generated before being randomly applied to the population.
\item Stratified randomization is the sorting of the population into groups by a certain characteristic before randomizing into either control or experimental groups.
\item Treatment allocation should be concealed until the moment of treatment initiation to avoid potential bias or manipulation.
\item Analysis of results should occur based on the initial randomization, irrespective of what may subsequently actually have happened to the subject “intention to treat analysis.”
\item Per protocol analysis is analyzing only those patients who strictly followed the assigned treatment protocol.
\item Most clinical trials use intention to treat to avoid bias and reflect real-world outcomes.
\end{itemize}
\begin{center}\small
\renewcommand{\arraystretch}{1.15}
\captionof{table}{The source contrasts analysis by assigned group with analysis restricted to protocol adherence.}\labtab{appendix-table-d2-intention-to-treat-and-per-protocol-analysis}
\begin{tabular}{p{0.45\linewidth}p{0.45\linewidth}}\toprule
\textbf{Intention-to-treat analysis} & \textbf{Per-protocol analysis} \\\midrule
Analyzes participants in their assigned groups.\par Preserves the benefits of random allocation. & Analyzes participants according to treatment received or adherence.\par May be biased when adherence differs between groups. \\\bottomrule
\end{tabular}
\end{center}
\subsubsection{Blinding}
\begin{itemize}
\item Blinding can reduce performance and measurement bias by keeping participants, clinicians, or assessors unaware of treatment allocation where feasible.
\item Blinding can reduce the influence of treatment expectations, including placebo and nocebo effects.
\item In a single blind experiment, only participants are blind to the treatment. In a double blind experiment, both participants and researchers are unaware of treatment allocation.
\item A placebo lacks the active treatment component. Expectations and treatment context may nevertheless affect reported symptoms and other outcomes.
\item The placebo should be identical to the treatment to maintain blinding.
\item The ‘double dummy’ placebo refers to both groups getting a different placebo similar to the alternate treatment regimen.
\end{itemize}
\subsubsection{Outcome Measurement}
\begin{itemize}
\item Each group must receive similar follow-up protocols.
\item The outcome chosen to be assessed in each group should be similar and appropriate to the research question.
\item The outcome should also be capable of being accurately and precisely measured.
\item Composite outcomes are outcomes which are a mix of multiple individual outcomes. They help assess the overall effect of a treatment.
\item Composite outcomes can improve efficiency, but their components should be clinically relevant and interpreted individually because an overall effect may be driven by a less important component.
\item Treatment effects may be summarized using absolute risk reduction (control event risk minus treatment event risk over a stated period), relative risk, or number needed to treat. For proportions, the NNT is the reciprocal of the absolute risk reduction and describes an average effect, not a guarantee of benefit for an individual.
\end{itemize}
\subsection{Prognostic Research}
\begin{itemize}
\item Prognostic research estimates or foresees risks for future conditions. It allows the prediction of the natural/future history of the disease.
\item It may be the basis for interventions to improve patient outcomes.
\item Prognostic research types:
\begin{itemize}
\item Natural history studies describe how a disease progresses over time without treatment.
\item Prognostic factor studies identify variables that influence disease outcomes.
\item Prognostic models combine factors into a risk score.
\item Predictive factor studies predict the response of a disease to treatment. 
\end{itemize}
\item The theoretical design of prognostic research can either be cause-effect (determining the effect of a factor, and answering “what will happen?”) or descriptive (eg. indicators for a disease or outcome).
\item In prognostic research, all variables potentially associated with the outcome, not necessarily causally, can be considered in the study. In other words, something associated with a disease may have not caused the disease.
\item The objective of prognostic research is to evaluate factors that might be modifiable to induce better patient outcomes.
\item The study sample should be the population at risk of developing the outcome of interest (having a particular condition).
\item The gold standard study design for prognostic research is the cohort study. It can be either prospective or retrospective.
\item Another less commonly used study design is the case-control, as it cannot determine absolute risk, and is prone to recall bias.
\item  Prognostic factors can be studied within RCTs as secondary analyses.
\item Candidate predictors can be obtained from patient demographics, clinical history, physical examination, disease characteristics, test results, and previous treatment. 
\item Predictors should be clearly defined, standardized, and reproducible. They should include only predictors that will be available at time of use.
\item Prognostic research focuses on outcomes that are relevant to patients, such as occurrence or remission, death, complications, tumor growth, or quality of life
\item Outcomes should be measured without knowledge of the predictors under study to prevent potential bias.
\end{itemize}
\section{Study Validity}
\begin{itemize}
\item Validity is about the accuracy of a measure.
\item Examination of the association between outcome of interest and exposure or intervention requires accurate and representative measures of the target population.
\item The research process:
\end{itemize}
\begin{enumerate}
\item Problem identification.
\item Literature review.
\item Research question and hypothesis forming.
\item Study design and planning.
\item Data collection.
\item Data analysis and interpretation of the results.
\item  Report writing.
\end{enumerate}
\begin{itemize}
\item An appropriate research question, well-defined study samples, and precise measurement methods enhances the internal validity of the study by minimizing biases and ensuring accurate causal inferences.
\item There are three main steps of study design: theoretical design (what correlation is expected) , data collection design and data analysis design.
\end{itemize}
\begin{itemize}
\item \textbf{Theoretical design}
\begin{itemize}
\item Research questions: What questions will the study address?
\item Background and significance: Why are these questions important?
\end{itemize}
\item \textbf{Data collection}
\begin{itemize}
\item Design: What is the time frame and epidemiologic design; how is the study structured?
\item Subjects: What are the selection criteria and sampling design; who are the subjects and how will they be selected?
\item Variables: Which predictor (independent), outcome (dependent) and confounding variables will be measured?
\end{itemize}
\item \textbf{Data analysis}: Which statistical issues, hypotheses, sample-size considerations and analytic approach will be used; how large is the study and how will it be analyzed?
\end{itemize}
\begin{center}\includegraphics[width=0.82\textwidth,height=0.68\textheight,keepaspectratio]{\reviewnotesimagepath d2-figure-research-design-components-analysis.png}\par\captionof{figure}{The schematic separates conceptual planning, participant and variable selection, and statistical analysis into linked design components.}\labfig{appendix-figure-d2-research-design-components-and-analysis}\end{center}
\begin{center}\small
\renewcommand{\arraystretch}{1.15}
\captionof{table}{Common study designs classified by the unit on which observations are made.}\labtab{appendix-table-d2-study-designs-and-observation-units}
\begin{tabular}{p{0.64\linewidth}p{0.24\linewidth}}\toprule
\textbf{Study design} & \textbf{Observation units} \\\midrule
\multicolumn{2}{l}{\textbf{Experimental studies}} \\
Randomized controlled trials & Individuals \\
Intervention trials & Individuals \\
Interventions in populations or communities & Groups \\\midrule
\multicolumn{2}{l}{\textbf{Observational studies}} \\
Ecological (correlation) studies & Groups \\
Case--control studies & Individuals \\
Cross-sectional studies & Individuals \\
Cohort studies & Individuals \\\bottomrule
\end{tabular}
\end{center}
\subsection{The Sampling Process}
\begin{itemize}
\item The ideal sampling process for studies is filtering the population for those meeting eligibility criteria and subsequently sampling them for the study, to allocate into groups as appropriate.
\item Representative sample populations result in more relevant observations.
\item Valid selection methods are by using randomized or categorized randomized sampling. This is often referred to as probability sampling.
\item Methods for probability sampling include:
\begin{itemize}
\item Simple random sampling.
\item Systematic random sampling (eg. every n people).
\item Cluster sampling (choosing the whole of a representative cluster, eg. a school).
\item Stratified sampling, which is sampling randomly from different categories that are established.
\end{itemize}
\item Convenience or nonrandom sampling often induces experimental errors.
\item Nonrandom sampling methods include:
\begin{itemize}
\item Convenience sampling, which is sampling the most accessible population to the researcher.
\item Purposive sampling is when the researcher uses their best judgement to select a sample most useful to their research purposes.
\item Snowball sampling recruits participants through recommendations from an existing participant.
\item Quota sampling divides the population into units and starts sampling from each unit until a quota is reached.
\end{itemize}
\end{itemize}
\subsection{Sample Allocation}
\begin{center}\includegraphics[width=0.82\textwidth,height=0.68\textheight,keepaspectratio]{\reviewnotesimagepath d2-figure-random-allocation-study-groups-outcomes.png}\par\captionof{figure}{A study sample is randomized into intervention and control groups, with outcomes measured after allocation.}\labfig{appendix-figure-d2-random-allocation-study-groups-outcomes}\end{center}
\begin{itemize}
\item Random allocation describes the process by which participants are put into the treatment and control groups of a study. There are four main types of randomization. For simple randomization, block randomization and stratified randomization, see above.
\item Covariate adaptive randomization uses the method of minimization by assessing the imbalance of sample size among several covariates. This method allows for the examination of previous participant group assignments to make a case-by-case decision on group assignment for each individual who enrolls in the study.
\begin{itemize}
\item Add the marginal totals of already assigned random samples to find the lower marginal total (the marginal total is the sum of any given row and column applicable to the participant).
\item Add the participant to the group with the lowest marginal total.
\end{itemize}
\item The probability of a false negative depends on sample size, effect size, variability, significance level, and study design.
\item A power analysis can be used to calculate the minimum sample size required so that one can be reasonably certain to detect an effect.
\item As the power increases, the chances of a false negative (Type II error) occurring decreases.
\item A confidence interval (CI) describes uncertainty around an estimate; its width reflects precision, but does not establish that the estimate is accurate or free from bias.
\item A wide confidence interval may reflect a small sample, high variability, or design features that reduce precision.
\item A small sample size does not mean that the results are wrong but rather that the data is consistent with a wide range of possible hypotheses.
\item A narrow or small confidence interval indicates a high precision.
\item Descriptive studies commonly estimate means or proportions with confidence intervals; hypothesis tests or power calculations may also be appropriate when the study question calls for them.
\item A confidence interval is a measure of the precision of a sample estimate. The width of a confidence interval depends on the sample size.
\end{itemize}
\begin{center}\small
\renewcommand{\arraystretch}{1.15}
\captionof{table}{The four possible decisions in a significance test, with typical design values of \(\alpha=0.05\), \(\beta=0.20\), and power \(0.80\).}\labtab{appendix-table-d2-hypothesis-testing-decisions-and-errors}
\begin{tabular}{p{0.25\linewidth}p{0.29\linewidth}p{0.29\linewidth}}\toprule
\textbf{Test decision} & \textbf{Null hypothesis true} & \textbf{Null hypothesis false} \\\midrule
Null hypothesis not rejected & Correct decision; probability \(1-\alpha\) & Type II error \(\beta\); probability \(1-\text{power}\) \\
Null hypothesis rejected & Type I error \(\alpha\); probability equals the significance level & Correct decision; probability equals the test's power \\\bottomrule
\end{tabular}
\end{center}
\subsection{Sample Size}
\begin{itemize}
\item When estimating sample size for descriptive studies, the investigator specifies the desired level and width of the confidence interval. The sample size can then be determined from the tables or formulas.
\item To estimate the sample size for that confidence interval:
\end{itemize}
\begin{enumerate}
\item Estimate the expected proportion with the variable of interest / not having it in the population (use the smaller proportion).
\item Specify the desired precision (total width) of the confidence interval.
\item Select the confidence level for the interval (e.g., 95\%).
\end{enumerate}
\begin{itemize}
\item For an analytic study, the sample should be estimated as follows:
\end{itemize}
\begin{enumerate}
\item State the null hypothesis and either a one- or two-sided alternative hypothesis.
\item Select the appropriate statistical test based on the type of determinant variable and outcome variable in those hypotheses (table below).
\item Choose a reasonable effect size.
\item Set \(\alpha\) and β. Specify a two-sided \(\alpha\) unless the alternative hypothesis is clearly one-sided.
\item Use the appropriate formula in an online calculator, or a statistics package to estimate the sample size.
\end{enumerate}
For a proportion with a 95\% confidence interval of width \(2d\), the estimated number of observations is
\[
n = \left(\frac{1.96}{d}\right)^2 p(1-p),
\]
where \(p\) is the expected proportion of the outcome, estimated from a previous, similar or pilot study, and \(d\) is the maximum tolerated error chosen by the investigator.
\subsection{Errors in Research}
\begin{itemize}
\item Erroneous inferences can be addressed through study design, implementation and analysis.
\item The goal is to maximize the validity of inferences from what was observed in the study sample to what is happening in the population.
\item There are three main sources of random errors in measurements: observer, instrument and subject (between subjects and between times).
\end{itemize}
Precision is assessed as the reproducibility of repeated measurements. Within-observer reproducibility compares measurements made repeatedly by the same person; between-observer reproducibility compares measurements made by different people.
\begin{center}\includegraphics[width=0.82\textwidth,height=0.68\textheight,keepaspectratio]{\reviewnotesimagepath d2-figure-measurement-repeatability-between-observers.png}\par\captionof{figure}{Repeatability is assessed by comparing repeated measurements made by one observer and by comparing measurements across observers.}\labfig{appendix-figure-d2-measurement-repeatability-within-between-observers}\end{center}
\begin{itemize}
\item There are two main approaches to minimizing random error and increasing the precision of measurements:
\begin{itemize}
\item Meticulous methodology increases accuracy:
\item[]
\begin{itemize}
\item Standardized measurements.
\item Certified observers.
\item Automation of instruments.
\item Unobtrusive measurements which subjects are unaware of.
\item Instrument calibration.
\item Blinding participants and observers, where feasible, can reduce some measurement and performance biases.
\end{itemize}
\item Standardization, calibration, and blinding can reduce particular measurement errors or biases, but none guarantees accuracy or eliminates every source of bias.
\item Larger samples generally reduce sampling variability in estimates, but do not eliminate random measurement error in individual observations.
\end{itemize}
\item The Hawthorne effect is when people change their behaviour due to observation.
\end{itemize}
\end{document}